\section*{ID8979: Definition of µ0µr(H+M)}
\paragraph{Metadata.}
PiID: 6818360900323 | RTEID: RTE:P34867.4510:T0.9022 | UUID: dc260609-a160-5608-8d27-dc060e47028e

\paragraph{Canonical Statement.}
Toroidal Curvature: Defines the minor and major radii \$(r,R)\$ to ensure the lemniscatic flux loop is perfectly centered within the vacuum chamber. [ID 422] HIGH-MU FERRITE CORE COMPOSITION \& SATURATION LIMITS Canonical Statutory Formula B = \$µ0µr(H+M) =⇒ dω\$ \$dt\$ = \$−γ(ω\$ − \$ω0)−K sin(2(θ\$ − \$θ0))\$ (22) Formal Technical Dissertation ID 422 defines the electromagnetic and dynamical properties of the rotating magnetic steering mechanism. The core utilizes a high-permeability (High-Mu) Manganese-Zinc (MnZn) ferrite, selected for its low eddy-current loss at the 17.2 kHz carrier frequency. Unlike standard rotors, this core operates under the "Pi-Model

\paragraph{Canonical Equation.}
\[
µ0µr(H+M) =⇒ dω
\]

\paragraph{Dictionary.}
Toroidal Curvature: Defines the minor and major radii (r,R) to ensure the lemniscatic flux loop is perfectly centered within the vacuum chamber.

\paragraph{Encyclopedia.}
Definition of µ0µr(H+M) is an algebraic definition, written µ0µr(H+M) =⇒ dω. It is an algebraic definition expressing the quantity directly in terms of the variables on its right-hand side. Within the Trawin Zero Point registry it serves as one element of the framework's mathematical narrative.
